\section{Limitations and Future Work}\label{sec:limitationsAndFutureWork}

\subsection{Reference counting overhead in signal combinators}\label{section:limitationsAndFutureWork:too_many_instructions}

As we described in Section~\ref{section:refcount:beans:counting}, the reactive semantics limit reuse in signal combinators to non-recursive combinators such as $stop$. In recursive combinators, the \textbf{reset} instructions are still inserted even though they will never succeed and these then force the insertion of additional \textbf{inc} instructions. 
The reason is that the reference counting algorithms always insert \textbf{reset}/\textbf{reuse} pairs whenever they encounter a suitable constructor expression, even if the memory will be needed again later.
A natural way to fix this would be to not insert \textbf{reset}/\textbf{reuse} in branches where the tail bound by the match is used in a recursive call. Preventing inserts at the level of branches would still allow reuse in combinators like $stop$ that have non-recursive and recursive branches. 

We are not the first to observe that the eager \textbf{reset}/\textbf{reuse} insertion can cause problems. 
\citeauthor*{frame_limited_reuse}~\cite{frame_limited_reuse} 
present a program similar to our $map$ combinator, but their concern is instead that the resulting code can hold onto memory for too long.
Their solution is to reverse the order in which \textbf{reset}/\textbf{reuse} instructions and \textbf{inc}/\textbf{dec} are inserted, so that \textbf{reset} instructions are only inserted when they can replace a \textbf{dec}.
This would also solve our issue of redundant and always failing \textbf{reset} instructions. 

\subsection{Specialised implementations of primitive combinators}

While all signal combinators can be implemented in RizzoLang, some of these may allocate more memory than necessary.
As an example consider the definition of $mk\_sig$ from Section~\ref{section:frp:rizzo}, 
If we unfold the $\triangleright$ operator, its definition becomes:
\[
\begin{aligned}
    &mk\_sig : \texttt{Later }a \rightarrow \texttt{Later }(\texttt{Signal }a) \\
    &mk\_sig \; d = (\lambda x. x :: mk\_sig \; d) \triangleright d \\
    \Updownarrow \\
    &mk\_sig \; d = \mathsf{laterapp} \; (\mathsf{delay} \; (\lambda a. \, a :: mk\_sig \; d)) \; d
\end{aligned}
\]
From this unfolded definition, we see that every (recursive) call allocates a fresh \textbf{ctor} value for the $\mathsf{laterapp}$, a \textbf{ctor} value for $\mathsf{delay}$, and a closure to captures $d$ for the lambda expression.
Given the shape of the later values generated by $mk\_sig$ never change, advancing such a later just reconstructs a new value that is structurally identical to the old one. This is wasted work.
% An efficient implementation would create the values on the first execution and reuse them for any future calls to $mk\_sig$. To do so we could define the function directly in the runtime, where it makes it possible to avoid these repeated allocations.
We suggest making $mk\_sig$ a primitive constructor for the \texttt{Later} type and adding a specialised implementation to handle such values in the advance semantics. This would allow us to avoid these repeated allocations. 
A similar optimisation could be made for the $\triangleright$ operator which is used in most combinators.

\subsection{The later applications forgotten to time}\label{sec:later_applications_forget}

Later applications, values constructed using $\mathsf{laterapp}$, can call arbitrary functions and are therefore not restricted by the operational semantics to return the same value each time. As a result, if $n$ signals depend on the same later application, that application will be evaluated $n$ times in a single time step, once for each dependency. This is inefficient both in terms of time and memory, as later applications may be expensive to compute, and they may also allocate signals on the signal heap. In the latter case, evaluating the same later application $n$ times will create $n$ identical copies of any signal allocated by the function. This duplication will cause excessive computation for as along as the program depends on those signals.

In Section~\ref{section:eval:heap_behaviour}, the signals $4$ and $6$ arise from exactly this situation: the same later application is evaluated twice, once for each dependency. Ideally, the signal $4$ would have been cached and returned for the second dependency, which would have prevented the creation of signal $6$ and the associated memory allocation. 

A current workaround is to turn the affected later computation into a signal and then apply the $\mathsf{tail}$ constructor to it:
\[
\mathsf{tail} \; (\texttt{""} :: mk\_sig \; (\mathsf{wait} \; console))
\]
This works because, according to the reactive semantics of $\mathsf{tail}$, advancing a value $\mathsf{tail} \; s$ returns exactly $s$.
This effectively means that the result of the later application is cached in the signal that we create as part of the workaround.
% For us to use $\mathsf{tail}$ we need to turn the \texttt{Later (Signal String)} from $mk\_sig$ into a \texttt{Signal String}. To do this, we take the \texttt{LaterApp} from \texttt{mk\_sig}, and apply it to the signal constructor (\texttt{::}) with a dummy head value (in this case the empty string). This instantiates a signal in the current time step, which we can use the \texttt{tail} constructor on to effectively cache the result of the later application. 
The downside is that we abandon the lazy signal creation. If the signal never produces a value, the allocation is unnecessary.
However, when the signal is used multiple times, we avoid repeated evaluation of the same later application and the associated memory allocations, which can significantly improvement performance and reduce memory usage.

A proper fix would be to make $\mathsf{laterapp}$ values behave like lazy values, so that the result is cached after the first evaluation.
However, the result may only be cached until a future time step when the $\mathsf{laterapp}$ produces a new value.
One way to implement this is to extend the representation of $\mathsf{laterapp}$ values with two additional fields. One field for the cached value and one for the time step for which the cache is valid.
%and one for a flag to indicate which time step the cached value is valid for.
Evaluation would then first check whether the cached value is valid. If so, then the cached value is returned. Otherwise, a new value is computed, stored in the cache and then returned.
%if the cached value is valid for the current time step and return it if it is, and otherwise compute the new value, store it in the cache, and return it.

% A similar problem can be observed for the Delay types, which are also just functions. However, the Delay type is less likely to be used in a way that causes multiple dependencies on the same Delay value, as it is more common to use the Delay type for delaying a signal by one time step, which would not cause multiple dependencies on the same Delay value within the same time step. 
